% Complete DRAM/storage research draft. Evaluation curves are analytical.
% Compile this directory independently with pdfLaTeX and BibTeX.
\documentclass[sigplan,10pt,anonymous,nonacm]{acmart}
\setkeys{acmart.cls}{balance=false}
\usepackage{subcaption}
\usepackage{tabularx}
\setcopyright{none}
\settopmatter{printacmref=false,printccs=false,printfolios=true}
\acmConference[EuroSys '27]{European Conference on Computer Systems}
 {April 19--23, 2027}{Rabat, Morocco}
\acmYear{2027}
\acmDOI{}
\acmISBN{}
\title{AegisKV: Workflow-Aware KV Checkpointing\texorpdfstring{\\}{ }Across DRAM and Storage}
\setlength{\emergencystretch}{1.5em}
\renewcommand{\bibfont}{\normalfont\normalsize}
\newcounter{designalgorithm}
\captionsetup[subfigure]{font=normalsize,labelfont=normalfont,textfont=normalfont}
\AtBeginDocument{\typeout{AEGISKV-GEOMETRY: width=\the\textwidth;
 height=\the\textheight; columnsep=\the\columnsep; baseline=\the\baselineskip}}
\begin{document}
\fancypagestyle{aegispreview}{\fancyhf{}\renewcommand{\headrulewidth}{0pt}
 \renewcommand{\footrulewidth}{0pt}\fancyfoot[C]{\normalsize\thepage}}
\pagestyle{aegispreview}
\begin{abstract}
Multi-agent applications accumulate reusable context while executing
interdependent LLM requests. A worker failure can lose this state and delay
several branches at once, yet protecting every context in fast memory is
expensive. We propose AegisKV, a workflow-aware manager for KV checkpoints
in shared DRAM and SSD or shared storage. AegisKV values DRAM retention,
storage demotion, and final eviction separately, combining normal reuse
with the workflow cost of recovery. It evaluates requests affected by a
common failure together, uses sparse demand and layer overlap to value
fast-tier residency, and screens online competitors using tier-local rent.
Versioned transitions preserve recovery validity during migration and
explicit fallback. The design requires no dedicated GPU backup region and leaves
the model's attention semantics unchanged. We connect trace-derived
workflow observations and existing reconstruction microbenchmarks to a
reproducible analytical study and a complete experimental protocol.
In 36 declared scenarios, the placement model improves on LRU in 17 and
ties in 19; these are model outcomes, not measurements of an implemented
system. The study identifies when storage bandwidth, prediction error,
and commit lag limit the benefit. End-to-end fault injection and normal
service measurements remain necessary to establish the proposed system's
performance.
\end{abstract}

\maketitle
\thispagestyle{aegispreview}
\raggedbottom

\section{Introduction}
\label{sec:introduction}

A multi-agent application completes when its required results are ready:
a tested patch, a report, or an answer assembled from parallel searches.
An LLM request's value therefore depends on its workflow position. A
critical-branch delay postpones completion; another branch may absorb it
before a join. Application-aware serving exploits this structure during
normal execution~\cite{lin2024parrot,pan2025kvflow,ni2026topas}. Failures
raise a further question: which histories should remain quickly
recoverable when one worker interrupts several requests together?

Long contexts make that choice expensive. Agents accumulate reasoning,
code, and tool results; losing their GPU-resident KV can require substantial
reconstruction. Checkpointing avoids that work~\cite{strati2024dejavu,jayakody2026ghostserve},
and load-aware recovery distributes interrupted requests~\cite{cao2026lumen}.
Hierarchical caches provide DRAM and storage capacity~\cite{qin2025mooncake,xie2026strata},
but keeping every history in the faster tier is costly.

Recency and recovery latency alone do not determine placement. A long,
recent history may belong to a branch that tolerates storage access,
while a smaller history lies on the critical path. High slack does not
justify deleting reusable history: subsequent normal calls may repeatedly
reconstruct it. Common-worker failures also couple requests through the
DAG and shared recovery resources. Accelerating only one of two delayed
branches may leave their join unchanged. Finally, demotion is asynchronous:
freeing DRAM before its storage replacement commits invalidates protection.

AegisKV manages checkpoints across shared DRAM and SSD or shared storage.
Its Tracker estimates remaining workflow time; the Master compares
complete plans under normal reuse and common-domain failures. Causal
sparse demand and layer overlap guide fast-tier retention, while tier-local
trial rent screens online competitors. The Router selects compatible
survivors; versioned manifests, reservations, and reader leases preserve
recovery validity. No dedicated Peer-HBM backup region is required.

Our contributions are a workflow objective that separates DRAM demotion
from final-copy deletion; a bounded online controller and safe transition
protocol; and reproducible analytical scenarios with a matched physical
experiment plan. The model exposes gains, ties, prediction sensitivity,
and regret against a feasible oracle. These results motivate the proposed
system; they do not measure its implementation. The central empirical
hypothesis is that workflow-aware protection improves finite-resource
allocation after accounting for healthy-service and migration costs.

% Background and Motivation: two background topics, two motivations,
% followed by a standalone failure-position example introducing AegisKV.
\section{Background and Motivation}
\label{sec:background-motivation}

\subsection{Background}
\label{sec:motivation-background}

\paragraph{Multi-agent workflows.}
Multi-agent workflows are becoming an important form of LLM application,
including coding assistants and research systems~\cite{wu2023autogen,anthropic2025multiagent}.
An agent may have a predefined role, such as a planner or reviewer, or be
instantiated as a subagent to complete a particular subtask. Each agent
may issue several LLM requests, invoke tools, and incorporate their results
before producing its output. These interactions form a workflow with
sequential dependencies, parallel branches, and joins. Multiple requests
using the same role may share a compatible prompt prefix, while successive
tool interactions extend an individual agent's history. Thus, a request,
an agent, and the complete user workflow describe different units of
execution and completion.

\paragraph{Failures at serving scale.}
LLM services distribute requests across many GPU workers, with each model
replica potentially spanning several devices. As deployments grow,
failures become an operational concern even when individual devices are
reliable: an unavailable device can interrupt an entire replica and shift
its work to the surviving fleet~\cite{cao2026lumen}. The interruption loses
both serving capacity and access to the failed worker's GPU-resident KV
cache. Recovery therefore requires more than rerouting a request; it must
also supply the state needed to resume computation. KV replication and
checkpointing systems already address this state-recovery
problem~\cite{strati2024dejavu,jayakody2026ghostserve,cao2026lumen}.
We focus on fail-stop GPU/worker failures with a compatible surviving
replica, retained token history and control state, and surviving backup
storage. Under this scope, the application-level consequences of the
interruption depend on both its workflow position and its recovery cost.

\begin{figure*}[t]
  \centering
  \captionsetup[subfigure]{font=normalsize,labelfont=normalfont,
    textfont=normalfont,justification=centering,singlelinecheck=false}
  \begin{subfigure}[t]{0.32\textwidth}
    \vspace{0pt}
    \begin{minipage}[c][2.25in][c]{\linewidth}
      \centering
      \includegraphics[width=\linewidth]{figures/workflow.pdf}
    \end{minipage}
    \caption{Four-branch workflow}
    \label{fig:motivation-workflow}
  \end{subfigure}\hfill
  \begin{subfigure}[t]{0.32\textwidth}
    \vspace{0pt}
    \begin{minipage}[c][2.25in][c]{\linewidth}
      \centering
      \includegraphics[width=\linewidth]{figures/workflow-gaps.pdf}
    \end{minipage}
    \caption{Workflow-maximum gaps}
    \label{fig:motivation-branch-gap}
  \end{subfigure}\hfill
  \begin{subfigure}[t]{0.32\textwidth}
    \vspace{0pt}
    \begin{minipage}[c][2.25in][c]{\linewidth}
      \centering
      \includegraphics[width=\linewidth]{figures/recovery-cost.pdf}
    \end{minipage}
    \caption{Context and Prefill cost}
    \label{fig:motivation-recompute}
  \end{subfigure}
  \caption{Motivating workflow-aware KV protection.
    (a) Selected Claude Code trace example with four parallel branches;
    A0/A5 denote Main before/after them. Bars preserve observed timing;
    dashed tails assume that all branch outputs are required at the join.
    (b) ECDF of the maximum branch-duration gap in each of 121 workflows
    with temporally inferred fan-in; it is not a slack distribution.
    (c) Measured Prefill latency for Qwen2.5-7B-Instruct-1M on two
    A100 PCIe 40\,GB GPUs (TP2, BF16, batch one), with MInference/DCA
    sparse Prefill above 32K, including the final one-token EXTEND.
    Reuse restores the same $(C-1)$-token KV and computes one new token.
    DRAM and shared-filesystem reads/H2D complete before computation.
    Shared FS is the measured DPC mount, with direct I/O and no resident
    Linux file pages; internal storage caches are uncontrolled.
    Full prefill starts from empty KV with 8K chunks.
    HBM/DRAM show medians and min--max ranges of three timings;
    Shared FS/full prefill each have one timing per length.
    K~=~1,024; M~=~$10^6$.}

  \label{fig:motivation}
  \Description{Left: a four-branch workflow with A2 completing last at
    151.6 seconds. Middle: an ECDF of workflow-maximum branch-duration
    differences with a median of 138 seconds and a 90th percentile of
    1173 seconds. Right: four measured Prefill curves for Qwen2.5-7B-Instruct-1M
    from 8K to 1,010,000 tokens: blue circles for HBM reuse, orange squares
    for serialized DRAM reuse, purple diamonds for serialized shared-filesystem
    reuse, and red triangles for full prefill. Both axes are logarithmic.}
\end{figure*}

\subsection{Motivation}

\paragraph{Workflow completion requires different protection priorities.}
\label{sec:motivation-workflow}
The user waits for the complete task, so we use workflow completion time
(WCT), from task release to completion of all required results, as the
primary objective. Minimizing an individual request's or agent's completion
time alone does not capture the effect of its dependencies.
Figure~\ref{fig:motivation}(\subref{fig:motivation-workflow}) illustrates
this distinction with a four-branch coding workflow. A2 completes last,
151.6\,s after the fork, whereas A1 completes at 60.4\,s. If the following
Main stage needs all four outputs, accelerating A1 alone does not advance
the join, while delaying A2 postpones downstream work. The same distinction
matters for fault protection: branches differ in how much additional delay
they can absorb. We call this tolerance \emph{slack}.

To characterize this heterogeneity beyond one example, we analyze the
public SemiAnalysis CC/Weka dataset collected from Claude Code coding
workflows~\cite{semianalysis2026cctraces}. Its filtered release contains
393 sessions. We retain overlapping subagent groups whose branches
complete and are followed by a Main request after the last branch finishes;
Main may also execute while the branches run. This yields 286 groups in
121 workflows. For each group, we compute the longest minus the shortest
branch duration; each workflow contributes its maximum such difference.
Figure~\ref{fig:motivation}(\subref{fig:motivation-branch-gap}) shows a
median of 138\,s and a 90th percentile of 1,173\,s. These statistics show
branch-duration heterogeneity among workflow maxima. The trace lacks
explicit result-consumption edges, so fan-in is inferred from timing and
the gaps are not a distribution of slack at failure time. The example and
distribution together motivate assigning protection priorities according
to workflow impact, rather than treating all agent requests equally.

\paragraph{Growing contexts make KV reconstruction increasingly costly.}
\label{sec:motivation-long-context}
An agent's later requests often include earlier reasoning, code, tool
outputs, and other agents' results. Without context reduction, this history
grows across interactions. KV caching avoids repeatedly processing that
history during normal execution. If a failure makes the active KV
unavailable and no usable copy survives, recovery must rebuild the missing
state before continuing; longer histories can therefore turn a brief
interruption into substantial additional computation.

Figure~\ref{fig:motivation}(\subref{fig:motivation-recompute}) measures
Qwen2.5-7B-Instruct-1M~\cite{qwen2025qwen25oneM} on two A100 PCIe
40\,GB GPUs (TP2, BF16 weights/KV, batch one). All paths share the
SGLang 0.5.10rc0 MInference vertical/slash configuration~\cite{jiang2024minference}
and DCA positions. Above 32K, including the final one-token EXTEND,
attention is sparse; shorter contexts use grouped dense DCA. Shared GQA
K/V and bounded query-head grouping avoid duplicating historical KV.

For total context $C=M+1$, full prefill processes $M$ tokens in chunks
of at most 8,192, then one separate token. Reuse restores that same
sparse-generated history and computes the same final token. HBM starts
resident; DRAM copies pinned host buffers to the GPU; Shared FS first
reads per-rank files from DPC into those buffers. File I/O, H2D and
computation are serialized, without compression or selective loading.
At 1,010,000 tokens the history occupies 57.92\,GB across both ranks.
Timing ends at synchronized model-body completion with all $C$ tokens'
KV resident, before vocabulary projection. Model loading and backup
creation are excluded; no failure is injected.

At that length, HBM, DRAM, Shared FS and full prefill take 1.305, 2.663,
10.156 and 571.21\,s. Reuse avoids historical computation; HBM also avoids
restore waits. The comparison does not quantify pipelined recovery.
Direct I/O and page-residency checks exclude Linux file-page hits;
DPC server/internal caches remain uncontrolled. This is neither cold-device
nor identified-SSD evidence.

The fixed synthetic text and sample counts are in the caption. HBM/DRAM
have one warmup; Shared FS has no read warmup. Short forwards and seven
operator checks precede timing; only 64K additionally has a complete
original all-head sparse reference, and 256K/1.01M have no full-prefill
warmup. Complete-history hashes, final hidden states, all-layer new KV
and greedy output match across paths; CUDA events verify serialization.
This validates the matched sparse paths, not dense-attention equivalence,
task quality or WCT. HBM is an execution-state reference, not Peer backup.

\begin{figure}[t]
  \centering
  \includegraphics[width=\linewidth]{figures/recovery-exposure.pdf}
  \caption{Workflow position changes the impact of recovery.
    Four independent single-branch failure counterfactuals for
    Figure~\ref{fig:motivation}(\subref{fig:motivation-workflow}),
    at 29.5\,s after the fork. Hatching is delay absorbed by join slack;
    solid fill is WCT loss. Reconstruction uses the active request's
    64-token hash-block input proxy and an assumed 64-token tail, with
    Qwen3-14B/A100 BF16 static estimates (batch and TP one).
    Other execution and queueing are held fixed; this is not a
    simultaneous four-branch fault experiment.}
  \label{fig:motivation-recovery-exposure}
  \Description{For A1 through A4, reconstruction costs are 35.2, 32.2,
    36.7, and 38.9 seconds, while corresponding WCT losses are 0, 32.2,
    17.6, and 17.0 seconds. A1 absorbs all reconstruction delay;
    A2 has zero slack and exposes all of its delay.}
\end{figure}

\paragraph{Failure position determines the workflow benefit of protection.}
The two motivations must be considered together: expensive reconstruction
does not imply an equally large WCT penalty. Figure~\ref{fig:motivation-recovery-exposure}
places an independent hypothetical failure on each active branch of
Figure~\ref{fig:motivation}(\subref{fig:motivation-workflow}) at the same
29.5\,s point. Under the stated join assumption, the observed completion
times provide offline slack, and static estimates provide the additional
reconstruction delay. A1's 35.2\,s cost fits within its 91.2\,s slack,
so it adds no WCT in this counterfactual. A2 costs only 32.2\,s to rebuild
but has zero slack, exposing the entire delay. A3 and A4 expose 17.6 and
17.0\,s, respectively. Hence, the branch with the cheapest reconstruction
can be the most important to protect. This controlled example illustrates
dependency effects; concurrent failures and shared recovery resources can
also delay other requests and must be evaluated separately.

\label{sec:motivation-goals}
These observations lead to AegisKV: a workflow-aware KV protection system
that allocates limited recovery state according to its effect on WCT.
The design connects request-level delay tolerance to the recovery value
of that request's KV, while accounting for DRAM, storage capacity, and shared I/O
consumed by protection. Existing valid cache copies remain usable for
recovery; additional protection must justify its cost to normal serving.
The next section develops the corresponding protection, placement, and
recovery mechanisms.

\section{AegisKV Design}
\label{sec:design}

AegisKV manages reusable KV checkpoints across shared DRAM and a larger
SSD or shared-storage tier. GPU HBM holds ordinary execution state and
bounded recovery workspace. The protected failure domain is a fail-stop
compute worker; the directory, token log, and advertised backing copies
must survive that failure. A host-local SSD does not protect against loss
of the same host. The following is a design specification; the new
mechanism figures are constructed examples or analytical schedules, not
measurements of an implemented controller.

\subsection{Design Insights}
\label{sec:design-insights}

\paragraph{Prioritize workflow exposure, while preserving normal reuse.}
A faster restore advances a workflow only when its delay exceeds the
available dependency slack (Figure~\ref{fig:motivation-recovery-exposure}).
However, a large-slack agent can reuse its history after a tool call.
Assigning that agent permanently to recomputation would waste serving
capacity. AegisKV therefore separates \emph{preserving a reusable copy}
from \emph{retaining it in the faster tier}. Shared immutable objects
serve both purposes without duplicating their physical byte charge.

\paragraph{Sparse access identifies candidates, not a complete backup.}
Native sparse attention can use only a subset of historical KV for a
query, but future selections remain uncertain. Figure~\ref{fig:design-insights}
shows that causal access history predicts future block demand unevenly
across models. At a 40\% block budget, past-frequency coverage is 77.8\%
for NOSA, 53.6\% for MiniCPM4, and 43.0\% for the experimental Seer path.
These observations motivate selective fast-tier retention and a mandatory
miss path, rather than discarding the remaining checkpoint.

\begin{figure}[t]
 \centering
 \includegraphics[width=\columnwidth]{figures/retry-kv-prediction.pdf}
 \caption{Historical access predicts distinct future KV unevenly.
 Three native PG19 captures/model use 32K history: queries
 32,257--32,768 predict blocks needed by 32,769--32,832.
 Objects are layer/64-token blocks containing all KV heads;
 mandatory blocks count against the 40\% budget. Seer uses an
 experimental Qwen3-8B Prefill gate~\cite{song2026seerqwen3}, threshold
 $10^{-3}$, with final-block dense override disabled.
 Bars/whiskers are document means/ranges, not recovery runs.}
 \label{fig:design-insights}
 \Description{Held-out native block coverage varies across three models;
 weak effects remain visible.}
\end{figure}

\paragraph{Overlap determines the marginal value of fast memory.}
DRAM-resident KV avoids the SSD read but still requires H2D transfer.
Early-layer reads can delay the start of replay, whereas later reads may
overlap with preceding replay computation. Their values also depend on
the current checkpoint length, uncommitted tail, demand reveal time,
resource contention, and workspace. A fixed resident layer count cannot
represent all requests. AegisKV evaluates feasible layer/block plans and
uses online rent to screen competing DRAM claims.

\subsection{System Overview}
\label{sec:system-overview}

Figure~\ref{fig:system-overview} separates workflow information, control,
and KV movement. The \emph{Tracker} estimates remaining durations and
slack; the \emph{Router} dispatches requests and selects surviving recovery
targets. The \emph{KV Master} values candidate placements and maintains
rental credits. Engine \emph{Adapters} expose sealed state and resource
pressure; \emph{Connectors} execute transfers. A surviving directory
records immutable objects, location generations, and checkpoint manifests.
Shared backing lets a compatible healthy DP restore state without a
preassigned peer-memory target.

\begin{figure}[t]
 \centering
 \includegraphics[width=\columnwidth]{figures/system-overview.pdf}
 \caption{AegisKV architecture. Both DPs use ordinary execution HBM.
 Solid blue arrows carry KV; dashed arrows carry control or metadata;
 gray workflow edges express dependencies. Shared DRAM retains valuable
 KV and SSD/storage preserves capacity. The surviving directory resolves
 recovery reads and publishes only validated copies.}
 \label{fig:system-overview}
 \Description{A single-column architecture connects an agent DAG,
 Tracker, KV Master, Router and directory to two symmetric DPs and
 shared DRAM and storage.}
\end{figure}

A physical object contains one layer's KV shard for a sealed token block.
Identity includes model/weight revision, prefix and position semantics,
layer, TP shard, format, and token interval. A legal placement group
includes the objects and selector state required by the native engine;
logical sparse selections, physical storage blocks, and kernel tiles need
not coincide. The engine chooses the physical block size; the 64-token
unit in Figure~\ref{fig:design-insights} is an experimental convention.

Prefill offers newly sealed groups. During decode, a configured,
block-aligned interval $Z$ triggers backup review. Arrivals, tool handoffs,
material demand/slack changes, and capacity pressure also trigger review.
A decision does not imply immediate commitment: if $M$ is the last
complete checkpoint frontier, the replay tail $N$ may exceed $Z$ when
copies lag. A checkpoint includes all historical dependencies needed to
resume native attention; DRAM and storage may jointly hold its objects.

\subsection{Workflow Tracker: Remaining Time and Slack}
\label{sec:workflow-protection}

The Tracker builds a residual DAG of unfinished execution segments.
Fork/join boundaries split parent work so child execution is not counted
again as parent waiting. A GBDT predicts remaining segment duration from
causally available role/model, elapsed time, known context, and completed
calls/tools. Entry predictions mask unavailable fields; training splits
separate complete workflows and time periods. Published trajectory
prediction~\cite{zhang2026heddle} motivates the features but does not establish
wall-clock prediction accuracy for AegisKV.

Critical-path analysis~\cite{kelley1959cpm} computes earliest and latest
finishes; their nonnegative difference is slack $\sigma_r$ inherited by a
request's segment. Requests on the same segment do not each receive an
independent copy of that margin. A refresh costs $O(|V|+|E|)$ after
prediction. Missing dependencies or invalid predictions fall back to
request recovery cost and normal reuse with zero slack. A new DAG version
invalidates outstanding valuation snapshots.



\subsection{Recovery-aware Router}
\label{sec:recovery-router}

Normal dispatch balances load and compatible prefix locality, retaining
Agent affinity when useful. The CC/Weka trace~\cite{semianalysis2026cctraces}
shows 97.90\% potential same-agent token coverage, with 0.37\% additional
coverage from other agents; these are potential prefixes, not cache hits. A tool handoff drops
an execution reference but does not delete the cached history or withdraw
its recovery claim.



On a confirmed failure, the Router fences the old source epoch, resolves
the last complete manifest, and chooses a healthy, layout-compatible
replica with feasible workspace. Queuing and readmission remain part of
recovery cost. Before acknowledging output, a surviving ordered log
records request ID, token position/value, and source epoch; prompt and
accepted tool inputs are also recoverable. Replay suppresses duplicate
delivery and rejects stale callbacks. Exactly-once external tool side
effects remain the orchestrator's responsibility.

\subsection{KV Master: Sparse Recovery Value and Online Rent}
\label{sec:kv-master}

\subsubsection{Valuing retention and demotion}
\label{sec:placement-value}

The Master compares complete, achievable placement plans. For an isolated
interrupted segment, recovery time $R_r$ exposes delay
$\max(0,R_r-\sigma_r)$. A common worker failure can interrupt several
branches: if two required branches each finish 10\,s late, accelerating
only one to 2\,s leaves their join 10\,s late. Thus, isolated block scores
cannot be added to obtain workflow benefit.

For current plan $P$ and candidate $P'$, we evaluate
\begin{equation}
 \Delta J=A(P',P)+\sum_e p_e\sum_w
       \bigl[T_w(P',e)-T_w(P,e)\bigr].
 \label{eq:value}
\end{equation}
Episodes $e$ cover normal next use and named compute-domain failures over
a fixed horizon. The weights $p_e$ are declared probabilities or expected
episode counts, not raw failure rates. A failure includes every request
interrupted in that domain. Profiles estimate additional segment delays
from surviving coverage, tail replay, I/O and target queues; a max-plus
DAG pass yields $T_w$. The transition term $A$ accounts for foreground
delay from copying and temporary occupancy, excluding costs already in
the episode terms. All terms have units of time.

Normal reuse remains in this objective even when failure probability is
zero or slack is large. DRAM demotion is valued against its actual
storage-backed continuation; last-copy deletion is valued against a
remaining checkpoint or reconstruction. Every shared owner is considered,
while each physical replica consumes capacity once. Pending copies help
only in episodes for which the profile predicts successful commitment.
No future failure time, prompt, or query is available to the policy.

\subsubsection{Sparse candidates and layer overlap}
\label{sec:sparse-placement}

For each request and layer, a bounded access summary ranks historical
blocks and includes native mandatory blocks. Candidates comprise legal
individual groups and small layer-ordered bundles, limited by a configured
work budget. The retained count $K_{r,l}$ describes the actual committed
DRAM groups; it can change as context, access, and competitors change.
There is no common resident depth $L$ or fixed attention-selection budget.
A sparse adapter predicts data demand without changing the model's native
selection. A miss reads storage or follows the declared reconstruction
fallback, as Figure~\ref{fig:sparse-recovery} illustrates.

\begin{figure*}[t]
 \centering
 \includegraphics[width=\textwidth]{figures/sparse-recovery.pdf}
 \caption{Sparse retention preserves full recovery semantics.
 (a) A complete historical checkpoint spans DRAM and SSD; dashed tail
 blocks still require replay. (b) A constructed prediction selects
 B0/B3, whereas the native query needs B0/B2. B2 is fetched on demand;
 B3 is an unused prefetch. Blue and tan consistently denote physical
 DRAM and SSD residency (D/S); unselected entries are unfilled.}
 \label{fig:sparse-recovery}
 \Description{Three layers have a complete mixed-tier history and an
 uncommitted tail. A predicted-block mismatch triggers storage access.}
\end{figure*}

The profile schedules SSD reads, H2D copies, and layer computation with
finite buffers. DRAM retention removes only the SSD stage. Only remaining
replay computation can hide transfers; avoided historical prefill is not
available overlap time. Figure~\ref{fig:design-pipeline} shows why equal
bytes at different layers can have different value. Figure~\ref{fig:overlap-frontier}
shows diminishing returns within each fixed compute profile. Larger
compute cost hides more later I/O but can increase total recovery time.
The runtime evaluates the current $M$, a declared distribution or bound
for unknown tail $N$, and demand-reveal constraints. It does not infer a
universal relationship between $Z$ and resident depth.

\begin{figure*}[t]
 \centering
 \includegraphics[width=\textwidth]{figures/recovery-pipeline.pdf}
 \caption{Analytical recovery schedules with equal-sized demand groups
 across four layers: SSD read 4\,ms/group, H2D 1\,ms/group, and replay
 3\,ms/layer. All four groups are correctly predicted at time zero;
 separate transient workspace holds them, and the three resources are
 independent. Total recovery is 20\,ms with all groups on SSD, 19\,ms
 with layer 4 retained in DRAM, and 16\,ms with layer 1 retained.
 These illustrative rates are not hardware measurements.}
 \label{fig:design-pipeline}
 \Description{Three resource timelines compare all-storage, late-layer
 DRAM retention, and early-layer DRAM retention under the same byte budget.}
\end{figure*}

\begin{figure}[t]
 \centering
 \includegraphics[width=\columnwidth]{figures/overlap-frontier.pdf}
 \caption{Analytical marginal value of retaining early-layer groups.
 Each curve fixes replay cost and shares the assumptions of
 Figure~\ref{fig:design-pipeline}; its zero-retention point is its matched
 baseline. Flat regions show reads already hidden by replay. The
 retained depth is an example candidate family, not a system-wide quota.}
 \label{fig:overlap-frontier}
 \Description{Recovery latency saturates at different retained group
 counts for three layer compute costs.}
\end{figure}

To capture limited complementarity, the candidate budget also permits
pairs of groups on required branches exposed to the same fault. Every
candidate is evaluated with its whole victim set through
Equation~\ref{eq:value}. Bounded search can still miss beneficial bundles;
we claim neither optimal placement nor submodularity.

\subsubsection{Online admission, rental, and eviction}
\label{sec:placement-plan}

AegisKV adapts credit aging from Landlord~\cite{young2002landlord} to screen
competitors independently in each capacity tier. For a physical replica
$x$ with bytes $b_x$, that tier's rent level $\rho$ and expiration key
$H_x$ represent remaining credit:
\begin{equation}
 c_x=b_x\max(0,H_x-\rho).
 \label{eq:rent}
\end{equation}
Credits are estimated delay saved, and rent has units of delay per byte;
it is not elapsed wall time. For each owner, the local proxy adds the next-use probability times
saved restore wait to fault-episode probabilities times the decrease in
$\max(0,R_r-\sigma_r)$, comparing the same group before and after a legal tier removal.
For DRAM this normally means storage access; for a last storage copy it
means the remaining checkpoint plus reconstruction. Each physical replica
has one byte charge and one tier-local ledger.
It only ranks candidates:
shared failures and migration costs still require the joint plan check.
No classical competitive guarantee transfers to this non-additive setting.

Credit is divided into owner shares while bytes are charged once. On a
new causal demand/checkpoint event, the affected owner's share is renewed
from its current proxy. Accrued rent is first settled proportionally across
old shares; other owners are not renewed. Repeated events with the same
version do nothing. Value-only refreshes may reduce credit but do not
refund paid rent. Thus, one active agent cannot perpetually renew an
unrelated owner's claim through a shared physical object.

All committed replicas in a tier age with its $\rho$, including replicas
outside its safe victim set $Q$. Membership in $Q$ requires no outstanding
physical read lease or copy dependency and an executable continuation:
a committed or reservable lower copy for demotion, or an explicitly
revised recovery frontier for last-copy removal. A running request's
logical ownership does not permanently pin every physical location. Credit
expiration alone never releases memory (Figure~\ref{fig:rental-competition}).

\begin{figure*}[t]
 \centering
 \includegraphics[width=\textwidth]{figures/rental-competition.pdf}
 \caption{Constructed rental step with equal-size groups. Trial rent
 increases by two credit units per byte, reducing credits 6/2/5 to
 4/0/3. The leased group ages from one to zero but is excluded from the
 safe victim set. The eligible zero-credit group is only a proposed
 demotion: the complete replacement must pass the workflow-value and
 reservation checks before any committed ledger changes.}
 \label{fig:rental-competition}
 \Description{All DRAM groups pay rent, while a reader lease prevents
 eviction irrespective of its credit.}
\end{figure*}

For optional admission, the Master simulates rent advancement using a
private snapshot. An incoming group's key is its credit per byte plus
current rent. While space is insufficient, it considers the eligible
resident with the smallest expiration key. Only an incoming key strictly
above both that key and trial rent can displace it; the trial rent advances
to that expiration and the victim joins the proposed plan. Ties retain the
resident. If no feasible victim remains, the admission fails. A free-space
admission does not advance rent.

The complete candidate-plus-victim plan must have a favorable upper cost
bound in Equation~\ref{eq:value}, with a configured hysteresis margin.
Rejected previews leave rent and locations unchanged. Accepted previews
atomically recheck versions, reserve resources, commit the new rent, and
enqueue transfers. Actual allocation waits for actual released bytes. A plan that needs a
lower-tier victim includes that victim and its recovery consequences in
the same preview; a DRAM action cannot assume unlimited SSD space.
Once execution starts, a later copy failure triggers reconciliation;
it cannot undo already completed migrations by rolling back a preview.

Required capacity relief has no incoming bidder: bounded rental victim
plans compete by predicted loss per releasable byte, subject to safety.
Storage deletion first removes unowned or superseded state. Under further
pressure, the Master may explicitly withdraw an active checkpoint claim,
valuing the resulting reconstruction for every owner. Before deleting its
last replica, it publishes a fallback to a complete older manifest, or to
prompt and surviving token-log replay if no older checkpoint remains.
Readers of the old generation must drain. It cannot retain the old
recoverability claim after deletion. If no valid fallback or executable
victim exists, it rejects/coalesces new writes and preserves the current
frontier. Identical storage capacity and fallback rules apply to baselines;
the ranking policy is the ablation variable.

\par\medskip\noindent\begin{minipage}{\columnwidth}
\raggedright
 \refstepcounter{designalgorithm}\label{alg:placement}
 \textbf{Algorithm \thedesignalgorithm: Online retention with trial rent}
 \par\smallskip
 \begin{tabularx}{\columnwidth}{@{}r@{\quad}X@{}}
 \toprule
 1 & Snapshot DAG, epochs, copies, leases, and byte ledgers \\
 2 & Reuse existing objects; offer newly sealed dependencies \\
 3 & Update causal demand summaries and affected credit shares \\
 4 & Form bounded sparse/layer candidates and safe victims \\
 5 & For each candidate, privately simulate rental displacement \\
 6 & Evaluate the complete feasible plan through the workflow DAG \\
 7 & Reject uncertain or unprofitable optional plans \\
 8 & Recheck versions; reserve destinations and transfer space \\
 9 & Commit trial rent; enqueue the selected placement changes \\
 10 & On pressure, choose safe relief; block admission if none \\
 11 & On completion/failure, reconcile actual copies and capacity \\
 \bottomrule
 \end{tabularx}
 \par\smallskip\noindent\textit{Preview capacity is not free memory; prediction is not a copy.}
 \end{minipage}\par\medskip

Profiles carry model/layout, size, tail, topology, queue load, support
range, uncertainty, and expiry. Missing or stale profiles freeze optional
promotions and fall back to bounded LRU with the same safety protocol.
Heap maintenance can be logarithmic, but complete cost includes candidate
construction, transfer profiling, and DAG passes. The implementation must
report decision latency and the effect of its candidate-work budget.

\subsection{Adapter and Connector: Safe Layered Recovery}
\label{sec:adapter}

Adapters report actual tier capacity and execution workspace; policy
windows update budgets without imposing a common $L$. Connectors enforce
reservations and bounded transfer queues for foreground reuse, demanded
recovery, and checkpoint/writeback. Configured service shares and bounded
quanta limit interference; demanded misses precede speculative reads
within the recovery share. These rules do not promise zero serving cost.

\begin{figure*}[t]
 \centering
 \includegraphics[width=\textwidth]{figures/checkpoint-lifecycle.pdf}
 \caption{Checkpoint and location transitions. (a) Pending tokens
 129--192 leave frontier $M=128$ unchanged; only dependency-complete
 commitment advances it to 192. (b) Demotion copies to storage,
 validates and commits the new location, redirects new readers, drains
 old leases, and then frees DRAM. (c) Last-copy deletion first publishes a
 complete older checkpoint or prompt/token-log replay. Numbers are illustrative;
 the protocol does not assume a particular backend persistence guarantee.}
 \label{fig:checkpoint-lifecycle}
 \Description{A complete manifest frontier differs from an in-flight
 copy; storage commitment and reader drainage precede DRAM release; last-copy removal explicitly changes the recovery path.}
\end{figure*}

Objects move through pending, committed, and retiring generations
(Figure~\ref{fig:checkpoint-lifecycle}). Advancing a frontier requires every
layer, TP shard, selector state, token interval, and matching surviving
log prefix. It is not the largest individually copied token position.
DRAM demotion follows copy--verify--commit--redirect--drain--free. The
backend must meet the claimed fault-domain persistence contract; buffered
write completion alone does not imply durable storage. Copy failure keeps
the previous authoritative location and releases only failed reservations.
Duplicate or stale-epoch callbacks cannot advance coverage.

Recovery acquires a manifest and leases its locations. The Connector
prefetches predicted native groups in layer order into bounded transient
workspace, while the engine replays the acknowledged tail. On-demand
selection resolves misses before executing the affected operation. Dense
layers use full demand; unavailable sparse metadata is restored or
reconstructed before selection. TP ranks obey matching collectives, and
batched replay is used only when it preserves the engine's native
semantics. A failure during prefill also starts from the last complete frontier,
possibly zero. Speculation may improve timing but never defines correctness.




\section{Evaluation}
\label{sec:evaluation}

We ask whether workflow information improves finite-tier allocation,
where that benefit disappears, and what it costs during normal service.
This version contains \emph{analytical results and an experimental design}.
It does not claim completed end-to-end deployment or failure experiments.
Existing trace observations and worker microbenchmarks in
Section~\ref{sec:background-motivation} retain their distinct evidence levels.

\subsection{Methodology and Reproducible Model}
\label{sec:eval-method}

\paragraph{Metrics and controls.}
For workflow $w$ under policy $p$, define the paired failure penalty
$D_{w,p}=\mathrm{WCT}^{\mathrm{fault}}_{w,p}-\mathrm{WCT}^{\mathrm{healthy}}_{w,p}$.
Each policy has its own matched no-failure run. Comparing policies uses
$D_{w,b}-D_{w,p}$, where positive values favor $p$; absolute healthy and
faulty WCT are also reported. The observation cohort includes directly
interrupted workflows and unaffected workflows sharing recovery resources.
Normal-service metrics include TTFT, inter-token delay, SLO goodput,
GPU-seconds, and per-tier request- and token-hit rates. Transfer bytes,
checkpoint age, management CPU time, and failed admissions expose costs
hidden by WCT alone.

\paragraph{Analytical assumptions.}
Four branches have 4, 8, 12, and 16\,GiB checkpoint groups and remaining
times 12, 10, 5, and 1\,s. Their next-use weights are 0.15, 0.35, 0.60,
and 0.85, normalized within a next-use episode. We assume two surviving
targets, a 16\,GiB DRAM budget, and sufficient finite storage (160\,GiB).
DRAM, storage, shared-network, and per-target H2D rates are respectively
20, 4, 12, and 20\,GiB/s. A common 0.25\,s recovery floor and 0.4\,s tail
replay complete the base profile. These are declared generic parameters,
not measurements or calibration of a named GPU or SSD.

An event calendar schedules 0.25\,GiB chunks through a shared DRAM or
storage reader, a shared network, and target-local H2D engines. Equal
chunks enter in round-robin request order. Ready operations take the
earliest available resource, with stable submission order breaking ties.
Tail replay waits for the request's state and shares target compute.
Transport stages overlap; this model does not claim layer-wise
transfer--compute overlap. Recovered branch times propagate through the
join. Placement completes before the modeled fault, so commit lag and
migration are studied separately rather than silently credited as zero.

The scenario grid has three state-size factors (0.5, 1, 2), four slack
factors (0, 0.3, 1, 3), and three storage rates (1, 4, 12\,GiB/s): 36
configurations, not 36 independent physical repetitions. DRAM scales with
state size. The scalar placement score uses failure-episode probability
0.05 and next-use-episode probability 0.02 over a 60\,s horizon; no-event
probability is the remainder. These illustrative probabilities are not a
production failure-rate estimate. The \emph{AegisKV model} greedily admits
whole groups by marginal score reduction per byte. It evaluates the
static value component, not tier-local rent, sparse demand prediction,
layer-wise replay, or the full asynchronous online algorithm.

\paragraph{Baselines.}
Storage-only retains all state in the capacity tier. LRU fills DRAM in
a declared most-recent order, A3 to A0; reuse-only prioritizes the assumed
next-use probabilities. The feasible oracle enumerates all 16 group
subsets under the same byte budget and minimizes the same score with true
durations. It is an allocation oracle within this restricted group family,
not a globally optimal online system. All-DRAM is an explicitly
capacity-relaxed latency reference. Recompute is separately available in
the model and required in physical tests. Named research systems receive
no fabricated performance curves.

\subsection{Workflow Outcomes and Capacity}
\label{sec:eval-workflow}

\begin{figure*}[t]
 \centering
 \begin{subfigure}{\columnwidth}
  \includegraphics[width=\linewidth]{figures/workflow-outcomes.pdf}
  \caption{Distribution across 36 configurations}
  \label{fig:eval-wct}
 \end{subfigure}\hfill
 \begin{subfigure}{\columnwidth}
  \includegraphics[width=\linewidth]{figures/capacity-frontier.pdf}
  \caption{Finite DRAM, fixed 40\,GiB state}
  \label{fig:eval-capacity}
 \end{subfigure}
 \caption{Analytical placement outcomes. (a) Each point is a deterministic
 scenario from the declared grid; the ECDF is not workload uncertainty.
 (b) Whole checkpoint groups explain discrete capacity transitions.
 All policies use the same queues, byte budget, and surviving targets.
 No confidence intervals or physical repetitions are implied.}
 \label{fig:eval-outcomes}
 \Description{An empirical distribution over enumerated analytical
 configurations and a capacity sweep compare LRU, reuse, workflow value,
 and a capacity-feasible allocation oracle.}
\end{figure*}

Figure~\ref{fig:eval-outcomes} shows the model's workflow penalties.
The AegisKV model improves on LRU in 17 configurations and ties in 19.
Mean penalties are 9.518\,s and 10.892\,s, respectively; these grid means
do not estimate a production improvement. The oracle's mean is 9.386\,s.
The largest greedy regret in the weighted placement objective is
0.204\,s, demonstrating that marginal admission does not solve the
non-additive problem exactly. LRU and reuse-only coincide under the
chosen recency and reuse order; this is a property of these scenarios,
not equivalence of the policies.

The capacity sweep provides two essential controls. With zero DRAM,
policies have the same storage path. Once all 40\,GiB fits, retention
policies converge. Between these endpoints, group sizes, shared queues,
and join sensitivity determine the result. Physical experiments must
sweep the working-set/DRAM ratio, include a no-contention region, and
retain per-workflow negative differences instead of reporting only a
favorable aggregate.

\subsection{Storage, Commit Lag, and Prediction Error}

\begin{figure*}[t]
 \centering
 \begin{subfigure}{\columnwidth}
  \includegraphics[width=\linewidth]{figures/storage-bottleneck.pdf}
  \caption{Storage path sensitivity}
 \end{subfigure}\hfill
 \begin{subfigure}{\columnwidth}
  \includegraphics[width=\linewidth]{figures/checkpoint-lag.pdf}
  \caption{Checkpoint lag and review interval}
 \end{subfigure}
 \caption{Analytical recovery sensitivity. Left: all-DRAM relaxes the
 16\,GiB limit and is only a reference. Right: the ideal phase is enumerated
 over $N=0,\ldots,Z-1$, with 100 generated tokens/s and 2.5\,ms replay/token;
 an additional commit lag adds acknowledged tokens to replay. Other
 parameters are fixed. Neither bandwidth nor replay rate is measured here.}
 \label{fig:eval-storage}
 \Description{Storage bandwidth and checkpoint lag sweeps expose I/O
 bottlenecks and the difference between a review trigger and commitment.}
\end{figure*}

Faster storage narrows the value of DRAM residency until another shared
resource dominates (Figure~\ref{fig:eval-storage}). This is why the main
comparison requires a tuned, pipelined DRAM/storage baseline rather than
only empty-cache recomputation. A shared filesystem must be labeled by
its actual mount, cache state, and network path; a directory name cannot
establish local NVMe behavior. Both client and server caching can change
the apparent bandwidth.

Review interval $Z$ is not a recovery bound. Under the idealized uniform
phase law, the mean uncommitted tail is $(Z-1)/2$ tokens; copying backlog
adds further acknowledged tokens. The right panel explicitly varies that
lag. Actual experiments sample observed commit frontiers and fail at
identical logical progress, with the same Workflow/Agent IDs. They must
also include failures during initial prefill, pending copies, and demotion.



Prediction error can reverse the benefit. In the base scenario,
alternating $-25\%$ duration perturbations increase the model penalty
from 1.613 to 4.558\,s, exceeding LRU's 3.621\,s. Outcomes use true
durations; this sensitivity example is not a GBDT accuracy estimate. Physical validation uses held-out
workflows and causal landmarks, reports absolute seconds-scale errors,
false-positive slack, fallback frequency, and placement regret, and
compares zero-slack, historical, learned, and oracle predictions. Both
systematic bias and out-of-distribution workloads are necessary controls.

\subsection{Attribution, Normal Service, and Safe Migration}

\begin{figure*}[t]
 \centering
 \includegraphics[width=\textwidth]{figures/policy-ablation.pdf}
 \caption{Analytical ablations, averaged equally over 12 size/slack
 configurations at 4\,GiB/s storage bandwidth. No DAG makes all branches
 equally critical; no reuse removes the normal-use term. Fault penalty
 and next-use restore delay are separate outcomes. Bars have no
 uncertainty intervals because the configurations are enumerated,
 not experimental repeats.}
 \label{fig:eval-ablation}
 \Description{Policy ablations compare workflow penalty with normal
 next-use restoration cost, retaining the tradeoff between the two objectives.}
\end{figure*}

Figure~\ref{fig:eval-ablation} removes dependencies and normal-use value
separately. The required physical comparison holds cache capacities,
checkpoint intervals, chunk sizes, transport scheduling, and failure
handling fixed. It changes only the decision rule. A separate comparison
against ordinary hierarchical caching plus a strong load-aware recovery
policy tests whether the combined system improves on an existing stack.
The revised online controller uses tier-local trial rent followed by
joint marginal revaluation. Required ablations compare direct revaluation,
rent without renewal, rent without aging, and full rent under the same
score and migration rules. Sparse retention and layer overlap are ablated
separately: recency versus causal access ranking, and serial loading versus
bounded pipelining, with equal committed coverage and transient workspace.
Online runs also fill storage and force last-copy fallback. The static
placement model in these plots does not validate those mechanisms.

Normal service is a first-class outcome: a bandwidth share cannot
guarantee a latency-overhead bound, because queueing depends on total
utilization. Full-system runs must therefore report healthy
WCT and SLO goodput as well as injected-failure improvements, including
long-context bursts, repeated Agent calls, and finite storage exhaustion.
A fault-free pass runs each original DAG once, without manufactured reuse
from repeated DAGs or clipping its initial calls.



Two invariants follow from the design: capacity reports follow physical
release, and common-failure evaluation follows the DAG. Analytical
examples do not validate concurrent code. The fault
campaign injects failures before/after reserve, copy completion, commit,
redirect, and lease release, plus full storage, duplicate callbacks, and
stale epochs. It checks exact checkpoint bytes and model-state parity,
recovery completion, duplicate-output suppression, and bounded queues.
Correctness failures are reported alongside latency and never dropped.

\subsection{Physical Experiment Plan and Statistical Analysis}
\label{sec:eval-plan}



Control-cost measurements include candidate construction, profile queries,
DAG passes, locks, and the outcome lost by bounding candidate work. An
assumed cost per node visit cannot establish controller scalability.

The prototype plan uses pinned SGLang/Mooncake revisions, two
failure-separated replicas and surviving storage, then 4/8 replicas for
load redistribution. Model cells include Qwen2.5-7B-Instruct-1M, Qwen3-8B,
and MiniCPM4-8B with exact weights, kernels and layouts. Dense/sparse paths
are separate; 32K--1M contexts are tested only where supported.

Workloads are instrumented coding, fan-out QA~\cite{zhu2024fanoutqa}, and
controlled sequential/fork--join DAGs. CC/Weka supplies timing/locality
observations, not recoverable prompts or true dependency edges. Sweep
DRAM/state ratio, storage capacity/bandwidth, and arrival rates at
30/60/85\% of the common baseline's saturation point. Inject at
20/50/80/95\% of the same Agent's logical token progress. Common-source
failures interrupt all colocated requests; request-only failures are
separate cells.

Each primary cell starts with ten independently launched paired campaigns.
Within a campaign, policies share workflow IDs, prompt/tool data, arrivals,
and logical failure events; each has its own healthy mate, with randomized
run order. Average paired effects within campaign, then bootstrap campaigns.
Report campaigns, workflows, and requests separately; more seeds do not
create new workloads. Tail quantiles remain descriptive when support is
inadequate. Preserve timeouts and implementation failures, report completion
and restricted completion time at a common horizon, and retain negative
results. Complete-case latency alone is insufficient.

The artifact specifies per-experiment controls, raw logs, baselines,
falsification conditions and paired aggregation. Identical CSV schemas
render analytical and measured data; measured mode checks raw hashes and
statistical provenance. Substituting data preserves figure layout, but
numbers, captions and claims must be revised to match the evidence.

\section{Related Work}
\label{sec:related}

\paragraph{Fault-tolerant LLM serving.}
D\'ej\`aVu streams and replicates KV~\cite{strati2024dejavu}; GhostServe
protects it with host-memory erasure-coded parity~\cite{jayakody2026ghostserve}.
RaidServe addresses TP resilience and imbalance through cyclic placement,
hybrid attention and recovery routing~\cite{xu2026raidserve}. LUMEN
coordinates checkpoints, request redistribution and capacity restoration
under reduced load~\cite{cao2026lumen}. Placement comparisons must use the same strong recovery scheduler.

\paragraph{Hierarchical context caching.}
Mooncake pools KV and schedules around locality~\cite{qin2025mooncake};
its SSD offload spills state under memory pressure~\cite{mooncake2026ssdoffload}.
CachedAttention uses hierarchical caching, layer-wise preload and
asynchronous saving~\cite{gao2024cachedattention}. Strata addresses
fragmented I/O, delay hits and cache-aware scheduling~\cite{xie2026strata}.
HCache restores intermediate activations while balancing I/O and
computation~\cite{gao2025hcache}. Our proposed distinction is common-failure workflow value
with safe demotion and explicit final-copy fallback.

\paragraph{Workflow-aware serving and prediction.}
Parrot optimizes application-level execution~\cite{lin2024parrot}; KVFlow
uses future invocation distance for prefix eviction~\cite{pan2025kvflow}.
Pythia and Cortex exploit workflow semantics~\cite{yu2026pythia,pagonas2025cortex}.
TOPAS schedules retention and execution using remaining paths, reuse and
movement cost~\cite{ni2026topas}; adding a critical-path score to caching
is thus insufficient novelty. The intended distinction is recovery after
a common worker failure with stale checkpoints and finite-tier transitions.
Heddle predicts trajectories~\cite{zhang2026heddle}, which does not establish
AegisKV's remaining-wall-time accuracy.

\paragraph{Placement mechanisms.}
Critical-path analysis~\cite{kelley1959cpm} and cost/size-based caching~\cite{young2002landlord}
provide building blocks. AegisKV's values depend on the DAG, shared queues,
and replicas, so Landlord's competitive bound does not transfer. The
small-instance feasible oracle measures approximation loss; independent
ablations must establish the value of rental, sparse ranking and overlap.

\section{Conclusion}
\label{sec:conclusion}

AegisKV combines normal reuse and common-failure workflow cost to manage
finite DRAM and storage. Its online decisions distinguish demotion from
final-copy deletion and preserve committed recovery state. Analytical
scenarios identify potential benefits and limits. The protocol and
reusable figures support the physical correctness and performance
experiments still needed to establish system-level gains.

\paragraph{AI assistance.}
Generative AI assisted drafting, experiment planning, and artifact code;
authors retain responsibility for verification.

\bibliographystyle{ACM-Reference-Format}
\bibliography{references}
\end{document}
